%% kapa: exam=1 round=33 year=2022 period=2 form=A subject=appraisal_law subject_ko=감정평가관계법규 seq=4 number=4
%% source: Q-Net 감정평가사 기출문제 (cst00309) artlSeq=5209997 file=raw/1008/5209997/2022년도 제33회 감정평가사 1차시험 2교시(감정평가 관계법규, 회계학).pdf
%% conversion: kordoc 4.24.51 -> tools/gen1.py (auto); review=proofread-sonnet; flags=none
\begin{question}{4}
「국토의 계획 및 이용에 관한 법률」은 중앙도시계획위원회와 지방도시계획위원회의 심의를 거치지 아니하고 개발행위의 허가를 하는 경우를 규정하고 있다. 이에 해당하는 개발행위를 모두 고른 것은?
\begin{framedtext}
ㄱ. 다른 법률에 따라 도시계획위원회의 심의를 받는 구역에서 하는 개발행위 \par
ㄴ. 「산림자원의 조성 및 관리에 관한 법률」에 따른 산림사업을 위한 개발행위 \par
ㄷ. 「사방사업법」에 따른 사방사업을 위한 개발행위 \par
\end{framedtext}
\begin{choices}
  \item ㄱ
  \item ㄴ
  \item ㄱ, ㄷ
  \item ㄴ, ㄷ
  \item ㄱ, ㄴ, ㄷ
\end{choices}
\end{question}
\answer{5}
